\documentclass[a4paper,twoside,12pt]{article}
\usepackage[top=1.5cm,left=1.5cm,right=1.5cm,bottom=2cm]{geometry}

\usepackage[utf8]{inputenc}
\usepackage[T1]{fontenc}
\usepackage{lmodern}
\usepackage[brazilian]{babel}

\usepackage{graphicx}
\usepackage{subfig}
\usepackage{psfrag}
\usepackage{cite}
\usepackage[cmex10]{amsmath}
\usepackage{amssymb}
\usepackage{upgreek}
\usepackage{microtype}
\usepackage{titling}
\usepackage{courier}
\usepackage{url}
\usepackage{hyperref}
\hypersetup{
	colorlinks=true,
	linkcolor=blue,
	urlcolor=blue
}
\usepackage[makeroom]{cancel}

% Use the 'listings' package to add code and the 'color' package to generate the highlights
\usepackage{listings}
\usepackage{color}
\definecolor{grey}{rgb}{0.6,0.6,0.6}
\definecolor{codegreen}{rgb}{0,0.6,0}
\definecolor{codegray}{rgb}{0.5,0.5,0.5}
\definecolor{codepurple}{rgb}{0.58,0,0.82}
\definecolor{backcolour}{rgb}{0.95,0.95,0.92}
\lstset{language=Python,				% set it to Python language highlighting
	backgroundcolor=\color{backcolour},
	commentstyle=\color{codegreen},
	keywordstyle=\color{magenta},
	numberstyle=\tiny\color{codegray},
	stringstyle=\color{codepurple},
	basicstyle=\ttfamily\small,		% set size of fonts
	breaklines=true,				% set automatic line breaking
	linewidth=\textwidth,			% set size of code box
	showstringspaces=false}			% show spaces as underscores only inside strings

\title{\vspace{-2cm} Processamento Digital de Sinais e Aplicações em Acústica\\ Soluções para Tutorial 07 - Série de Fourier e a Transformada Discreta de Fourier}
\author{\url{https://github.com/fchirono/Aulas_PDS_Acustica}}

\begin{document}
	\date{}
	\maketitle

%\setcounter{section}{1}
\section{Condições de Correspondência DFT -- Série de Fourier}
\setcounter{subsection}{1}

\subsection*{Tarefa 1}

Para comparar a equação da DFT (Eq. 8) com a equação de análise da Série de Fourier na forma exponencial (Eq. 7), podemos aproximar a integral de análise em $t$ por uma somatória em $n$ (onde $t = n \cdot \Delta t$):

\begin{equation}
	\int_{t=0}^{T_0} f(t) \ dt \approx \sum_{n=0}^{N_t-1} f(n \cdot \Delta t) \ \Delta t.
\end{equation}

Continuando a derivação, encontram-se três condições devem ser satisfeitas simultaneamente para que esta aproximação seja \emph{exata}:

\begin{enumerate}
	\item \textbf{Normalização dos coeficientes}: $C_k = X[k] \ / \ N_t$
	
	\item \textbf{Periodicidade do sinal analisado}: o sinal deve ser periódico com período $T_0 = 1/f_0$.
	
	\item \textbf{Análise de um período inteiro}: a duração $T = N_t \cdot \Delta t$ do vetor de amostras deve ser exatamente igual ao período $T_0 = 1/f_0$ do sinal. Quando esta condição é satisfeita, $f_0 = \Delta f$ e o $k$-ésimo harmônico ($k \cdot f_0$) da frequência fundamental coincide exatamente com o $k$-ésimo \emph{bin} da DFT ($k \cdot \Delta f$).
\end{enumerate}

A violação da segunda condição leva ao \emph{vazamento espectral} (\textit{spectral leakage}), estudado em detalhes no Tutorial 08. Naturalmente, o critério de Nyquist ($f_s \geq 2 f_{max}$, onde $f_{max}$ é a maior frequência presente no sinal) ainda deve ser satisfeito.


\end{document}
